\section{Overview}
The rapid proliferation of digital communication channels has been matched by an equally rapid evolution in cybercrime methodologies. Traditional approaches to identifying unsolicited, malicious, or fraudulent electronic messages have transitioned from rudimentary rule-based filtering to sophisticated, high-dimensional machine learning frameworks. This chapter reviews the state-of-the-art literature relevant to the foundational pillars of the CyberShield project: Short Message Service (SMS) and text-based fraud detection, threshold uncertainty and active learning, and the integration of Human-in-the-Loop (HITL) methodologies in cybersecurity. This chapter is organized into four sections: classical ML approaches to spam detection (Section 2.2), NLP tooling (Section 2.3), Human-in-the-Loop and active learning paradigms (Section 2.4), and a gap analysis situating CyberShield within the existing literature (Section 2.5).

\section{Machine Learning in Spam and Fraud Detection}
The academic investigation into automated filtering of electronic messages conceptually began with email spam. However, due to the constrained character limits, informal syntactic structures, and specific colloquialisms prevalent in SMS and instant messaging platforms (like WhatsApp), traditional email detection algorithms perform poorly when directly transposed to short-text environments.

\subsection{Evolution of Text Vectorization and Modeling}
Almeida et al.~\cite{almeida2011contributions} introduced the SMS Spam Collection v.1, which remains one of the most widely used public benchmarks for SMS spam classification research. Their comparative study of standard algorithms - including Support Vector Machines (SVM) and Naive Bayes (NB) - established that Term Frequency–Inverse Document Frequency (TF-IDF) vectorization combined with linear classifiers produces competitive baseline metrics on short-text data. Subsequent work has demonstrated that adversarial obfuscation techniques (character substitution, homoglyph attacks, whitespace manipulation) can degrade detector performance, motivating the need for character-level features and dynamic retraining pipelines such as those employed in CyberShield.

Building upon this, Abid et al. \cite{abid2022spam} conducted a comprehensive comparative analysis of traditional machine learning against deep learning methodologies for SMS spam detection. Their research demonstrated that while Deep Learning models (like LSTMs and CNNs) achieve modest performance gains on highly complex, non-linear linguistic structures, traditional algorithms (specifically optimized SVM and Logistic Regression models) provide significantly faster inference times and lower computational overhead, making them well-suited for real-time deployment. Furthermore, tree-based and linear models provide inherent feature explainability, which is a critical requirement in forensic cybercrime investigations.

Similarly, Roy et al. \cite{roy2020deep} investigated deep learning mechanisms but emphasized the importance of extensive dataset curation. Their findings underscored the ``garbage in, garbage out" paradigm; highly imbalanced datasets result in models that collapse towards the majority class (typically legitimate messages). Consequently, optimal detection relies heavily on synthetic minority oversampling (SMOTE) or hard negative mining to penalize the misclassification of edge-case fraud vectors.

\subsection{Smishing and Phishing URL Detection}
A significant subset of text-based cybercrime relies on delivering malicious payloads via Uniform Resource Locators (URLs), broadly categorized as Smishing (SMS Phishing). Jain and Gupta \cite{jain2021smishing} proposed an integrated machine learning approach that does not solely rely on the semantic text of the message, but actively parses the embedded URLs to extract lexical features (e.g., URL length, presence of hyphens, highly entropic subdomains). Their work demonstrated that combining NLP semantic analysis with lexical URL analysis reduces the false positive rate.

Further solidifying this, Sahingoz et al. \cite{sahingoz2019machine} demonstrated that lexical URL features alone - domain length, character entropy, presence of suspicious tokens - can achieve over 97\% phishing detection accuracy on large-scale datasets. The consensus across multiple reviewed studies indicates that attackers frequently exploit URL shortening services (e.g., bit.ly, tinyurl) to obfuscate malicious domains from localized text scanners. To counter this, CyberShield incorporates a Velocity Intelligence module specifically designed to track the submission frequency of these shortened artifacts, flagging them as malicious based on submission volume regardless of the domain's reputation at the time of submission.

\section{Natural Language Processing Tooling}
The implementation of machine learning models requires robust frameworks capable of transforming raw, unstructured text into mathematical representations. The Natural Language Toolkit (NLTK) established the academic standard for programmatic linguistics. Their comprehensive documentation covers the necessity of stop-word removal, stemming, and lemmatization to reduce the dimensionality of the vector space, ensuring that algorithms are not distracted by syntactically necessary but semantically hollow conjugations.

Although contextual embedding approaches such as Word2Vec and BERT offer richer semantic representations, two practical constraints favored TF-IDF for this deployment: (i) the inference latency budget of 500ms per request, which contextual transformer models exceed on CPU-only infrastructure, and (ii) the requirement for token-level explainability via inspectable model coefficients, which is straightforward for linear models over TF-IDF features but opaque for transformer embeddings. For the predictive modeling phase, the \textit{scikit-learn} library has become the de facto standard in the machine learning community. Their architecture provides standardized, interchangeable modules for vectorization, enabling researchers to rapidly prototype ensemble mechanisms. CyberShield utilizes these exact theoretical mathematical constructs to perform its soft-voting probability aggregation.

\section{Human-in-the-Loop (HITL) and Active Learning}
A glaring deficiency in modern commercial fraud filters is their deterministic binary nature. Models are mathematically forced to classify a message as $1$ (Fraud) or $0$ (Legit), often making arbitrary decisions on messages that sit squarely on the algorithmic decision boundary. 

\subsection{The Necessity of the Human Operator}
Holzinger \cite{holzinger2016interactive} explored this paradox extensively in the context of health informatics, positing the question: \textit{When do we need the human-in-the-loop?} Holzinger argues that in domains where the cost of a false positive is catastrophic, fully autonomous ``black-box" AI is suboptimal. Instead, algorithms should be utilized to perform the heavy lifting of parsing vast datasets, extracting the ambiguous ``uncertain" samples, and presenting them to a human expert. The human leverages contextual domain knowledge that the AI lacks, annotates the ambiguous sample, and feeds it back into the training loop. This specific philosophy explicitly underpins CyberShield's multi-tier thresholding logic, where probabilities between $0.30$ and $0.70$ are deferred to admin authorities rather than guessed autonomously.

\subsection{Concept Drift and Active Learning Paradigms}
This mechanism of querying a human expert is formally known in academic literature as Active Learning. Settles \cite{settles2009active} provides the definitive survey on Active Learning, outlining various query strategies. The most relevant to threat intelligence is ``Uncertainty Sampling," where the model queries the human operator for the labels of instances about which it is least certain. Furthermore, active learning directly mitigates the phenomenon of \textit{Concept Drift} \cite{gama2014survey}, where the statistical properties of the target variable (e.g., spam patterns) change rapidly over time. Settles shows that selecting boundary cases for human labelling enables a machine learning model to reach a target accuracy with substantially fewer labelled instances than random passive sampling.

\subsection{Hard Negative Mining}
Addressing class imbalance and boundary ambiguity programmatically leads to the concept of Hard Negative Mining. Shrivastava, Gupta, and Girshick \cite{shrivastava2016training} presented an online hard example mining algorithm initially designed for region-based object detectors in computer vision. While originating in computer vision, the core principle of heavily penalizing high-confidence errors transfers directly across domains and is highly effective when adapted for NLP fraud detection. In the context of CyberShield, a ``hard negative" is a highly suspicious legitimate message (e.g., an urgent banking OTP) that the model incorrectly flags as a scam. By programmatically extracting these failures during the cross-validation phase and forcing the model to retrain explicitly on them, the decision boundary is shifted away from the legitimate class distribution, reducing the risk of false alarms on borderline cases.

\section{Gap Analysis and CyberShield Contribution}
The reviewed literature reveals a significant disconnect between theoretical machine learning accuracy and practical, real-world deployment in cybercrime mitigation. Most studies conclude upon achieving a marginally superior accuracy metric on static, sterile datasets. 

There is a notable gap in literature regarding end-to-end socio-technical pipelines. Academic models do not translate their raw probabilistic outputs into actionable intelligence for the victim (such as an automated NCRP FIR schema). Furthermore, public datasets suffer from extreme temporal decay; a model trained predominantly on 2011-era SMS spam vectors (such as the Almeida dataset~\cite{almeida2011contributions}) is likely to underperform against contemporary WhatsApp KYC and PAN-card update scams that did not exist when the dataset was compiled. 

CyberShield addresses this gap by combining the established robustness of TF-IDF ensemble modelling with Holzinger's Human-in-the-Loop paradigm~\cite{holzinger2016interactive} and iterative Hard Negative Mining~\cite{shrivastava2016training}. The result is an auto-calibrating fraud detection platform that integrates with national cybercrime reporting standards rather than producing only research-grade metrics.

\subsection{Comparison with Recent Systems}
To definitively contextualize the novelty of CyberShield, Table \ref{tab:lit_comparison} illustrates a direct architectural comparison between recent notable academic systems and our proposed architecture.

\begin{table}[H]
    \centering
    \resizebox{\textwidth}{!}{
    \begin{tabular}{|l|l|c|c|c|p{5cm}|}
    \hline
    \textbf{System / Authors} & \textbf{Algorithm Core} & \textbf{Artifact Neutralization} & \textbf{HITL Calibration} & \textbf{Uncertainty Handling} & \textbf{Primary Limitation / Gap} \tabularnewline \hline
    Almeida et al. (2011) \cite{almeida2011contributions} & SVM, Naive Bayes & No & No & None & \raggedright Binary only; high false positive rate on edge cases. \tabularnewline \hline
    Roy et al. (2020) \cite{roy2020deep} & CNN / LSTM & No & No & None & \raggedright Computationally heavy; lacks explainability for cyber intelligence. \tabularnewline \hline
    Jain \& Gupta (2021) \cite{jain2021smishing} & Lexical URL + NLP & Yes (URLs only) & No & None & \raggedright Does not extract Phone Numbers; no feedback loop for concept drift. \tabularnewline \hline
    Abid et al. (2022) \cite{abid2022spam} & Optimized LR / SVM & No & No & None & \raggedright Black-box predictions without automated legal reporting integration. \tabularnewline \hline
    \textbf{CyberShield (Current)} & \textbf{Ensemble (LR+SVM+NB)} & \textbf{Yes (URLs + Phones)} & \textbf{Yes (Active Learning)} & \textbf{4-tier Probabilistic} & \raggedright \textbf{Currently English/Hinglish only; no native multilingual support.} \tabularnewline \hline
    \end{tabular}
    }
    \caption{Architectural Comparison of CyberShield against contemporary spam detection systems.}
    \label{tab:lit_comparison}
\end{table}
